(1) 界面无自由电荷意味着法向位移连续，\(D_{1n}=D_{2n}=Q/A\)。因此在忽略边缘效应的均匀场近似下，\(E_1=Q/(A\varepsilon_1)\)、\(E_2=Q/(A\varepsilon_2)\)，两者同向但一般并不相等，较小介电常数层承受更强电场。
(2) 两层电压是 \(V_1=E_1d_1\)、\(V_2=E_2d_2\)，故 \(V=Q(d_1/\varepsilon_1+d_2/\varepsilon_2)/A\)。等效电容 \(C=Q/V=A/(d_1/\varepsilon_1+d_2/\varepsilon_2)\)，也满足 \(1/C=d_1/(\varepsilon_1A)+d_2/(\varepsilon_2A)\) 的串联关系。
(3) 第 \(i\) 层储能为 \(\varepsilon_iE_i^2Ad_i/2=Q^2d_i/(2A\varepsilon_i)\)。两层相加给 \(U=Q^2(d_1/\varepsilon_1+d_2/\varepsilon_2)/(2A)=Q^2/(2C)\)，与从等效电容计算的能量完全一致，且显然非负。
(4) 压降比 \(V_1/V_2=(d_1/\varepsilon_1)/(d_2/\varepsilon_2)\)。若 \(\varepsilon_1=\varepsilon_2=\varepsilon\)，电容退化为 \(\varepsilon A/(d_1+d_2)\)；若 \(d_2\to0\)，得到 \(\varepsilon_1A/d_1\)。界面无自由电荷保证法向位移连续，不应错误要求法向电场相等。
